\begin{table*}[!t]
\centering
\caption{Stage 1 regime, spike, and 1\% VaR diagnostics. Average QLIKE hides a common practical weakness: most forecasts remain too smooth on spike days, and tail coverage can fail even when average volatility loss is competitive.}
\label{tab:stage1-risk-diagnostics}
\scriptsize
\renewcommand{\arraystretch}{0.95}
\setlength{\tabcolsep}{2pt}
\resizebox{\textwidth}{!}{%
\begin{tabular}{@{}l r r r l L{0.28\textwidth}@{}}
\toprule
Model & Extreme-regime QLIKE & Severe spike underpred. & Spike pred./actual & 1\% VaR violations & Interpretation \\
\midrule
GARCH(1,1) & 6.758 & 86.7\% & 0.273 & 16/745 (\(p=0.00631\)) & Strong average benchmark, but strict tail coverage fails \\
AdvGARCH-BestQLIKE & 6.756 & 88.9\% & 0.251 & 15/745 (\(p=0.01453\)) & Primary advanced model; spike timing remains imperfect \\
HAR-Parkinson & 6.335 & 88.9\% & 0.245 & 16/745 (\(p=0.00631\)) & Competitive average QLIKE with similar tail undercoverage \\
EWMA(\(\lambda=0.90\)) & 8.345 & 88.9\% & 0.266 & 18/745 (\(p=0.00101\)) & Complete industry baseline, but weaker under extreme volatility \\
LSTM-LogTarget & 20.256 & 100.0\% & 0.056 & 56/726 (\(p=3.81{\times}10^{-31}\)) & Low smooth forecasts severely understate spikes and tails \\
Hybrid-LogTarget-Small & 19.929 & 100.0\% & 0.060 & 48/726 (\(p=5.02{\times}10^{-24}\)) & Low MAE does not translate into risk reliability \\
\bottomrule
\end{tabular}
}
\end{table*}
